\subsection{\problemblock{*Print} Data Block}\label{sec:trajectory-block-print}

The trajectory-level \problemblock{*print} block is similar to the
\problemblock{*file} block in \cref{sec:trajectory-block-file}, except that it
writes trajectory variables to the standard printout file instead of to a user-named
column file. The standard file has the same basename as the problem file and uses the
extension \texttt{.out} rather than \texttt{.prb}. It contains page breaks and
headings formatted for a printer with 132 characters per line and 60 lines per page.

One table of trajectory data is written for each \problemblock{*print} block. Any
number of blocks may appear in a trajectory, with a maximum of ten trajectory
variables in each block. For example,

\begin{lstlisting}[style=taosmanual]
*print time alt range mach gamgd dynprs
\end{lstlisting}

produces output like \cref{fig:trajectory-printout-example}.

\begin{figure}[htbp]
  \centering
  \begin{minipage}{0.96\linewidth}
    \lstinputlisting[style=taosmanual,basicstyle=\ttfamily\scriptsize]{examples/chapter04/trajectory-printout-example.txt}
  \end{minipage}
  \caption{Trajectory printout example.}
  \label{fig:trajectory-printout-example}
\end{figure}

The output begins with a page heading, title, and column headings, which are
repeated on each page, followed by the trajectory values. Units and formats are
controlled by \problemblock{*units/fmt}.

A line is printed at the beginning and end of every segment and at each print interval
within the segment. A blank line separates segments so they can be identified easily.
Radar and relative-vehicle variables, whose names begin with \texttt{rad} or
\texttt{rel}, require a square-bracket subscript. For example,
\statevar{radrng[2]} is range from station 2 to the current trajectory, and
\statevar{relvel[4]} is the closing velocity of trajectory 4 relative to the current
trajectory.
